\subsection{Plano de Testes}
\label{sec:plano_de_testes}

Para garantir a corretude, a robustez e a manutenibilidade da plataforma, foi adotada uma estratégia de testes automatizados em múltiplas camadas, inspirada na pirâmide de testes. Cada camada possui um escopo e um objetivo distintos, validando diferentes aspectos do sistema, desde a lógica de negócio isolada até o fluxo completo da aplicação em um ambiente simulado.

\subsubsection{Testes de Unidade (Unit Tests)}
Os testes de unidade formam a base da pirâmide de testes. Eles são rápidos, focados e executados com frequência durante o desenvolvimento.
\begin{itemize}
    \item \textbf{Objetivo:} Validar a lógica de negócio de cada componente individual da camada de aplicação (e.g., \textit{Command Handlers} e \textit{Queries}) em completo isolamento de suas dependências externas.
    \item \textbf{Implementação:} Utilizando o framework \texttt{pytest} e a biblioteca \texttt{unittest.mock}, todas as dependências de infraestrutura — como repositórios de banco de dados, o \textit{Event Bus} e serviços externos — são substituídas por \textit{mocks}. Isso permite testar a lógica de negócio pura de um handler (e.g., "dado o comando X, o handler Y deve chamar o método Z do repositório e publicar o evento W?") de forma determinística e instantânea.
\end{itemize}

\subsubsection{Testes de Integração (Integration Tests)}
Os testes de integração ocupam o meio da pirâmide e servem para validar a colaboração entre a lógica da aplicação e um componente de infraestrutura real, neste caso, o banco de dados.
\begin{itemize}
    \item \textbf{Objetivo:} Garantir que a camada de persistência de dados está funcionando corretamente. Isso inclui verificar se os modelos de domínio do SQLAlchemy estão corretamente mapeados para o esquema do banco de dados e se as consultas SQL geradas pelos \textit{Repositories} se comportam como o esperado.
    \item \textbf{Implementação:} Os testes de integração, localizados em \texttt{tests/integration}, utilizam um banco de dados PostgreSQL real, orquestrado via Docker. Uma \textit{fixture} do \texttt{pytest} (definida em \texttt{tests/integration/conftest.py}) é responsável por criar o esquema do banco de dados antes da execução dos testes e por limpá-lo ao final. Cada teste executa um \textit{Command Handler} contra um repositório real e, em seguida, utiliza a mesma sessão de banco de dados para verificar se os dados foram persistidos corretamente na tabela, validando a integração entre a aplicação e a camada de dados.
\end{itemize}

\subsubsection{Testes de Ponta a Ponta (End-to-End Tests)}
Os testes de ponta a ponta (E2E) estão no topo da pirâmide. Eles são os mais lentos e complexos, mas fornecem o maior nível de confiança, pois validam o sistema completo da perspectiva do usuário.
\begin{itemize}
    \item \textbf{Objetivo:} Validar que o pipeline de processamento de mídia assíncrono e orientado a eventos funciona corretamente de ponta a ponta, desde a requisição HTTP inicial até a verificação do resultado final.
    \item \textbf{Implementação:} Foi criado um ambiente de teste completo e isolado com o arquivo \texttt{docker-compose.test.yml}. Este ambiente orquestra todos os serviços da aplicação: a API, o worker Celery, o banco de dados de teste e o Redis. Um contêiner dedicado, o \texttt{test-runner}, executa o script de teste E2E (localizado em \texttt{tests/e2e}). O script utiliza a biblioteca \texttt{requests} para fazer chamadas HTTP reais à API, simulando o fluxo de um usuário: (1) autenticação, (2) upload de um arquivo de mídia, (3) "polling" do status da mídia até que o processamento seja concluído, e (4) verificação dos artefatos finais (transcrição e resumo), garantindo que todos os componentes do sistema distribuído colaboraram com sucesso.
\end{itemize}

Esta estratégia de testes em múltiplas camadas garante que, embora cada componente seja validado individualmente, a integração entre eles e o comportamento do sistema como um todo também sejam rigorosamente verificados, proporcionando alta confiança na corretude e na estabilidade da aplicação.
